% Section 3: Census Sampling Uncertainty (~1,600 words)

\subsection{The Nature of PL 94-171 Uncertainty}

The redistricting algorithm uses PL 94-171 population data \citep{pl94171}
as input. PL 94-171 is the Census Bureau's official redistricting data
file, released approximately 18 months after Census Day. It is based on
a complete enumeration attempt---not a survey sample---but ``complete''
is aspirational. The Census Bureau's own Coverage Measurement program
documents systematic undercounting and differential coverage rates by
race, ethnicity, and urbanicity \citep{brown2018demographic,
daponte2021census}.

For redistricting purposes, two aspects of census uncertainty matter:

\textbf{Overall undercount rate} affects the denominator used to compute
the ideal district population: $P^* = P_\text{total} / k_s$, where
$k_s$ is the number of districts apportioned to state $s$.
If $P_\text{total}$ is underestimated by $\epsilon\%$, the ideal
district size is underestimated by $\epsilon\%$, and the population
balance requirement is calibrated to a slightly wrong target.

\textbf{Differential undercount} by geography affects which tracts
appear to satisfy the population balance constraint. Tracts in
communities with high differential undercount (typically minority,
urban, and young-adult-heavy communities) will appear smaller than they
truly are, potentially affecting how the algorithm groups them.

\subsection{Published Undercount Estimates}

The Census Bureau's 2020 Post-Enumeration Survey (PES)
\citep{uscensusbureau2022net} estimated net coverage errors by state
and demographic group. Key findings:
\begin{itemize}
  \item National net undercount: $-0.24\%$ (undercount of total
    population by approximately 749,000 people)
  \item Range by state: $-1.1\%$ (Texas, estimated undercount)
    to $+0.8\%$ (Hawaii, estimated overcount)
  \item Differential undercount by race: Black ($-3.3\%$), Hispanic
    ($-4.99\%$), American Indian on reservations ($-5.64\%$) relative
    to non-Hispanic White ($+1.64\%$ overcount)
\end{itemize}

These estimates have their own sampling uncertainty, reported as 90\%
CIs by the Census Bureau. For the national figure: $-0.24\%$
(90\% CI: $[-0.56\%, +0.08\%]$). For Texas: $-1.1\%$
(90\% CI: $[-1.7\%, -0.5\%]$).

\subsection{Delta Method for Population Deviation}

The population deviation of a district from ideal is:
\begin{equation}
  \delta_i = \frac{P_i - P^*}{P^*},
  \quad P^* = \frac{P_\text{state}}{k_s}
\end{equation}
where $P_i$ is the population of district $i$ and $P^* $ is the
ideal district population. The redistricting algorithm requires
$|\delta_i| \leq 0.005$ (the $\pm 0.5\%$ constitutional standard).

Suppose the true population $P_i^{\text{true}} = P_i (1 + \epsilon_i)$
where $\epsilon_i$ is the local coverage error rate. Then the true
deviation from ideal is:
\begin{equation}
  \delta_i^{\text{true}}
    = \frac{P_i(1+\epsilon_i) - P^*(1+\bar\epsilon)}{P^*(1+\bar\epsilon)}
    \approx \delta_i + (\epsilon_i - \bar\epsilon) + O(\epsilon^2)
\end{equation}
where $\bar\epsilon = \sum_i P_i \epsilon_i / P_\text{state}$ is the
state-level mean coverage error rate.

By the delta method, the variance of $\delta_i^{\text{true}} - \delta_i$
is approximately:
\begin{equation}
  \text{Var}(\delta_i^{\text{true}} - \delta_i)
    \approx \text{Var}(\epsilon_i - \bar\epsilon)
    = \sigma^2_\epsilon (1 - 2 w_i)
\end{equation}
where $w_i = P_i / P_\text{state}$ is the district's share of state
population and $\sigma^2_\epsilon$ is the between-tract variance
of coverage errors.

\subsection{Empirical Calibration}

We calibrate $\sigma^2_\epsilon$ using the Census Bureau's
tract-level undercount estimates from the 2020 PES experimental
microdata \citep{uscensusbureau2022net}. Across all tracts, the
estimated SD of tract-level coverage error rates is
$\sigma_\epsilon \approx 0.015$ ($1.5\%$).

For a typical congressional district ($P_i / P_\text{state} \approx 1/k_s$),
the delta method gives:
\begin{equation}
  \text{SE}(\delta_i^{\text{true}} - \delta_i)
    \approx \sigma_\epsilon \sqrt{1 - 2/k_s}
    \approx 0.015 \text{ for large } k_s.
\end{equation}

The 95\% CI for the true deviation, given an observed deviation $\delta_i$,
is therefore approximately:
\begin{equation}
  \delta_i \pm 1.96 \times 0.015 = \delta_i \pm 0.029.
\end{equation}

This is a wide interval---wider than the $\pm 0.5\%$ constitutional
tolerance. However, this reflects \emph{uncertainty about the true
population}, not about whether the algorithm satisfies the constitutional
standard given the official census data. The algorithm's target is the
official population, and its deviation is measured against that target.
Census uncertainty is relevant for interpreting whether population
balance in the official data implies population balance in reality.

\subsection{Implications for the $\pm 0.5\%$ Standard}

The practical implication is:

\begin{proposition}
A district that satisfies $|\delta_i| \leq 0.005$ using official
PL 94-171 data may have a true population deviation up to
$|\delta_i^{\text{true}}| \approx 0.034$ due to differential
census undercount, with 95\% probability.
\end{proposition}

This does not invalidate algorithmic redistricting; no redistricting
method (algorithmic or legislative) can achieve true population equality
when the census itself has coverage errors. What it does mean is that
claims of ``equal population districts'' are legally precise (equal
as measured by official data) but not literally exact.

For the purposes of UQ reporting, we recommend:
\begin{itemize}
  \item Report population deviation as: $|\delta_i| \leq 0.005$
    (per official 2020 Census PL 94-171 data; true deviation may
    differ by up to $\pm 0.029$ at 95\% confidence due to census
    coverage uncertainty).

  \item In states with documented high undercount (Texas, Florida,
    Georgia) use state-specific undercount estimates from the 2020 PES
    to bound coverage error more precisely.

  \item Note that differential undercount ($-3.3\%$ for Black populations,
    $-4.99\%$ for Hispanic) may mean majority-minority districts are
    actually larger than recorded, potentially benefiting VRA compliance
    in practice.
\end{itemize}

\subsection{Impact on Ideal District Size}

The ideal district size for the 2020 apportionment is:
$P^* = 331,449,281 / 435 = 761,952$ persons (official).
The 90\% CI for the true total U.S. population, using the national
PES estimate $(-0.24\%, 90\%\text{ CI: }[-0.56\%, +0.08\%])$, is
$[331,449,281 \times (1 - 0.0056),\ 331,449,281 \times (1 + 0.0008)]
= [329,595,000,\ 332,714,000]$.
This translates to ideal district sizes in the range $[757,000,\ 764,000]$---
a variation of $\pm 3,500$ persons ($\pm 0.46\%$) in the target.

For states near the congressional apportionment margin (where a small
population change would change the number of districts), this uncertainty
matters. Texas ($k=38$) and Montana ($k=2$ per the 2020 cycle) were
close calls; the PES CI on national population is too wide to rule out
different apportionments under alternative census methodologies.
